\section{Homología relativa}

Concluimos el desarrollo de los aspectos que necesitamos de homología con la \textit{homología relativa}. En esencia, en homología relativa lo que hacemos es ignorar todas las cadenas de un subespacio. La construcción es similar a lo que ya conocemos\dots

\begin{definition}
    Sean $X$ un espacio y  $A\subset X$ un subespacio. Definimos $C_k(X,A)$ como el grupo cociente
    \begin{equation*}
        C_k(X)/C_k(A),
    \end{equation*} 
    luego las cadenas de $A$ se vuelven triviales.
\end{definition}
 
Ahora, para los operadores de borde vayamos por partes. En primer lugar, para que tenga sentido definirlos debe ocurrir que $\partial_k$ lleve $C_k(A)$ a $C_{k-1}(A)$. 

En el caso de un espacio triangulado $A$ debe ser \textit{subespacio triangulado}, es decir, un subespacio cubierto únicamente por poliedros completos de la triangulación original. Para CW-complejos el subespacio debe ser un \textit{subcomplejo}, es decir una colección de celdas completas del espacio.

Una vez tenemos esto, el operador de borde en homología relativa no es más que el operador de borde en homología seguido del paso al cociente. En efecto, gracias a que $\partial_k(C_k(A))$ está contenido en $C_{k-1}(A)$, estos inducen un operador de borde en el cociente dado por
\begin{equation*}
    \partial^{rel}_k = q_{k-1}\circ\partial_k,
\end{equation*} 
donde los $q_k$ forman la familia de funciones de paso al cociente. Omitiremos la etiqueta ``rel'' siempre que el contexto lo permita.

Observamos que en este caso también se verifica $\partial^{rel}_{k-1}\circ \partial^{rel}_k = 0$ ya que se anulan antes de pasar al cociente, luego los $q_k$ forman un morfismo de cadena.

Una vez tenemos los grupos de cadenas y los operadores de borde definidos, se construyen los \textit{grupos de homología relativa} de la misma manera que en homología celular y simplicial.

\begin{nota}
    Como cabría esperar, si el subespacio $A$ es por ejemplo contractible o vacío, se recuperan los grupos de homología no relativa.
\end{nota}
\bigskip

\begin{proposition}
    Para cualquier par $(X,A)$ con las propiedades necesarias según el contexto como establecido arriba se tiene la siguiente secuencia exacta
    \begin{align*}
        \cdots \longrightarrow H_n(A) \longrightarrow H_n(X) \longrightarrow H_n(X,A) \longrightarrow H_{n-1}(A) \longrightarrow H_{n-1}(X) \longrightarrow \cdots \longrightarrow H_0(X,A) \longrightarrow 0.
    \end{align*}
\end{proposition}

Demostrar esto es uno de los objetivos principales de la sección de homología relativa de \cite[p.~115]{hatcher2002} en la que nos hemos basado para esta sección.

Un ejemplo de la utilidad de esta secuencia exacta para computar los grupos de homología relativa es el siguiente.
{\color{red} lo del disco sin la frontera que hice con angel.}

\bigskip

%%%%%%%%%%%%%%%%%%%%%%%%%%%%%%%%%%%%%%%%%%%%%%%%
%%                                            %%
%%            NUMEROS DE BETTI                %%
%%                                            %%
%%%%%%%%%%%%%%%%%%%%%%%%%%%%%%%%%%%%%%%%%%%%%%%%
\section{Invariantes relacionados con la homología}

Introducimos brevemente dos invariantes topológicos relacionados con la homología de un espacio que serán centrales en la última sección.

\begin{definition}
    Sea $X$ un espacio topológico. Definimos el \textit{$k$-ésimo número de Betti de X} como 
    \begin{equation*}
        b_k(X) = \text{rank} H_k(X).
    \end{equation*}
\end{definition}


Para la definición general de rango de un grupo abeliano referimos al lector a la página $85$ de \cite{FuchsInfiniteAbelianGroupsI}. En el caso de que los grupos de homología sean finitamente generados, como es el de un espacio compacto, su rango queda determinado por su descomposición según el \textit{teorema de clasificación de grupos abelianos finitamente generados}. Así, $\text{rank} H_k(X)$ es $r$ en
\begin{equation*}
    H_k(X) = \underbrace{\mathbb{Z}\oplus\dots\oplus\mathbb{Z}}_{r} \oplus \mathbb{Z}_{d_1} \oplus\dots\oplus \mathbb{Z_{d_m}}.
\end{equation*}
\bigskip

\begin{definition}
    Sea $X$ un espacio topológico cuyos grupos de homología son finitamente generados y tal que $H_k(X)=0$ para todo $k>N$. La \textit{característica de Euler-Poincaré de X} es 
    \begin{equation*}
        \chi (X) = \sum^N_{k=0} (-1)^k b_k(X).
    \end{equation*}
\end{definition}

\begin{nota}
    Una definición habitual de la característica de Euler-Poincaré en términos de los poliedros de una triangulación se puede encontrar en la página $39$ de \cite{munoz2021}. En ella se pide que $X$ sea compacto para que tenga sentido sumar recorriendo los poliedros de la triangulación, aquí ponemos una condición más débil. 
    
    Dado que no siempre contamos en nuestro caso con una triangulación conviene tomar la versión más general de arriba en términos de los números de Betti.
\end{nota}

Tanto para los números de Betti como para la característica de Euler-Poincaré se pueden cambiar $b_k(X)$ y $\chi(X)$ por $b_k(X,A)$ y $\chi(X,A)$ y se conservan las definiciones para los grupos de homología relativa.

\bigskip